\section{Implementation Decisions and Trade-offs}

\subsection{Connection reuse vs. a fresh connection per packet}

The HTTP/2 generator opens one \texttt{httpx} client with HTTP/2 enabled and
keeps it open for the whole run, firing several requests at once over that
single connection to actually use multiplexing rather than just claiming to.
The QUIC generator does the same thing for the same reason: a QUIC handshake
costs at least one extra round trip, and if the generator reconnected for
every request that handshake cost would swamp the latency numbers and make
QUIC look artificially slow next to HTTP/2 and MQTT, both of which also keep a
connection open.

The TCP/UDP generator does the opposite on purpose. For every single TCP
packet it opens a new connection, sends, and closes it again. This was not an
oversight. A fresh connect/close cycle produces a real SYN, SYN-ACK and FIN
sequence in the capture, which is exactly what is needed later for the
Failure Visibility analysis and for telling raw TCP apart from the
application-level protocols in a protocol hierarchy view.

\subsection{Mode and pattern as two separate knobs}

The TCP/UDP generator splits "what does a packet look like" from "how is it
paced over time". \texttt{mode} controls size and base interval: normal mode
uses a fixed 512 bytes every 100 ms, stealth mode draws a random size between
64 and 1400 bytes with Poisson-distributed timing. \texttt{pattern}
independently controls the sending cadence on top of that: constant, periodic
burst, an extra random gap, or a linear ramp. Keeping these two orthogonal
made it possible to test, for example, stealth-mode sizing combined with a
burst cadence, which would not have been possible if mode and pattern were
the same setting. The other three generators (HTTP/2, QUIC, MQTT) do not have
a separate \texttt{mode} the way TCP/UDP does, but they share the same
\texttt{pattern} field and the same four options, so periodic bursts and
ramps are available on every protocol, not just the one that originally
needed them for behavioral fingerprinting.

\subsection{The Poisson pattern on every protocol, not just one}

The assignment only asks for a random, Poisson-distributed pattern somewhere
in the system. This implementation puts it on all four generators using
\texttt{random.expovariate(1/interval)} (or NumPy's exponential distribution
for the TCP/UDP generator), which produces exponentially distributed gaps
between sends with the same average rate as the constant pattern. The reason
to extend it everywhere was to be able to show one single Wireshark capture
where HTTP/2, QUIC and MQTT all use Poisson timing at once, alongside TCP/UDP
stealth mode, rather than only ever demonstrating it on one protocol.

\subsection{Warmup/cooldown and the per-protocol ramp pattern are two separate mechanisms}

An early version of this system tried to satisfy the assignment's "ramping
(linear increase over time)" sending pattern by reusing the warmup/cooldown
machinery: a controller-side function, \texttt{\_ramp\_rates()}, already
steps every generator's rate field linearly from zero up to its configured
target over the warmup window, and back down to zero over the cooldown
window, sending the new value through the normal PATCH endpoint at each
step. The gradual change is visible as a slope in a Wireshark I/O graph, so
at first glance it looked like one mechanism could cover both requirements.

On closer reading of the assignment, this was not a good fit. Ramping is
listed alongside constant, periodic burst, and random as one of several
sending-pattern \emph{options}, meant to be selected per protocol and per
phase, the same way a phase can put one protocol on \texttt{periodic\_burst}
while another stays \texttt{constant}. The warmup/cooldown ramp cannot do
that: it always ramps every generator's rate together, and only once, at the
very start and end of an entire profile run, never in the middle of it on a
single protocol.

Each generator was therefore given its own \texttt{ramp} pattern, alongside
\texttt{constant}, \texttt{periodic\_burst}, and \texttt{random}, configured
with \texttt{ramp\_start\_rate}, \texttt{ramp\_end\_rate}, and
\texttt{ramp\_duration}. The generator tracks a wall-clock timestamp for when
it last switched into ramp mode and computes its current effective rate by
linearly interpolating between the two rate values over the configured
duration, independently of the controller and of any other protocol running
at the same time. A phase can now, for example, set MQTT to ramp from 5 to
180 messages/sec over 40 seconds while HTTP/2 in the same phase stays at a
flat rate -- something the global warmup/cooldown ramp alone could never
express. The two mechanisms can coexist (warmup/cooldown still ramps the
generic rate field for every generator at the edges of a run; the per-protocol
ramp pattern runs inside whichever phase requests it), with one minor
interaction worth noting: while a generator is in \texttt{ramp} mode, it
reads its rate from \texttt{ramp\_start\_rate}/\texttt{ramp\_end\_rate}
rather than the generic \texttt{rate} field, so a warmup/cooldown PATCH
targeting that generic field has no visible effect on it until the phase's
pattern changes away from \texttt{ramp}. None of the bundled configuration
profiles trigger this, since their ramp phases all sit between the warmup
and cooldown windows, not overlapping them.

\subsection{Adaptive Control}

The controller can run an autonomous control loop, either because a YAML
phase turns it on or because it was switched on manually. Every few seconds it
checks each generator's error rate, and where reported its latency, against
configurable thresholds, and scales the rate up or down by a multiplier
bounded on both ends. This goes past the baseline requirements, but it ties
directly into the resilience patterns covered in the course: the system reacts
to an injected failure on its own, without anyone touching a slider.
